\subsection{Vector Spaces, Bases, and Matrices}
Linear algebra is the study of vector spaces and linear mappings. Vectors are, by definition, elements of some \textit{vector space} $\mathcal{V}$ and satisfy the usual rules of addition and scalar multiplication, e.g.,
\[
    \text{if } \v{x} \in \mathcal{V} \text{ and } \v{y} \in \mathcal{V}, \text{ then } \alpha \v{x} + \beta \v{y} \in \mathcal{V} \text{ for all } \alpha, \beta \in \R \text{ (or } \C\text{)}.
\]

We will be dealing mostly with vectors in the Euclidean vector space $\R^n$ for some $n$. That is, we view the points of $\R^n$ as objects that can be added up and multiplied by a scalar, e.g., $(x_1, x_2) + (y_1, y_2) = (x_1 + y_1, x_2 + y_2)$ for points in $\R^2$. Sometimes it is convenient to work with the complex vector space $\C^n$ instead of $\R^n$.

Vectors $\v{\nu}_1, \dots, \v{\nu}_k$ are called \textit{linearly independent} if none of them can be expressed as a linear combination of the others; that is, if $\alpha_1 \v{\nu}_1 + \dots + \alpha_n \v{\nu}_n = \v{0}$, then it must hold that $\alpha_i = 0$ for all $i = 1, \dots, n$.

\begin{defn}{Basis of a Vector Space}{BasisOfAVectorSpace}
A set of vectors $\mathcal{B} = \{\v{\nu}_1, \dots, \v{\nu}_n\}$ is called a \textit{basis} of the vector space $\mathcal{V}$ if every vector $\v{x} \in \mathcal{V}$ can be written as a unique linear combination of the vectors in $\mathcal{B}$.
\begin{equation*}
    \v{x} = \alpha_1 \v{\nu}_1 + \dots + \alpha_n \v{\nu}_n.
\end{equation*}
The (possibly infinite) number $n$ is called the \textit{dimension} of $\mathcal{V}$.
\end{defn}
Using a basis $\mathcal{B}$ of $\mathcal{V}$, we can thus represent each vector $\v{x} \in \mathcal{V}$ as a row or column of numbers
\begin{equation}
    [\alpha_1, \dots, \alpha_n] \qquad \text{or} \qquad
    \begin{bmatrix}
        \alpha_1 \\ \vdots \\ \alpha_n
    \end{bmatrix}.
    \label{eq:VectorRep}
\end{equation}
Typically, vectors in $\R^n$ are represented via the \textit{standard basis}, consisting of unit vectors (points) $\v{e}_1 = (1, 0, \dots, 0), \dots, \v{e}_n = (0, 0, \dots, 0, 1)$. As a consequence, any point $(x_1, \dots, x_n) \in \R^n$ can be represented, using the standard basis, as a row or column vector of the form \cref{eq:VectorRep} above, with $\alpha_i = x_i, i = 1, \dots, n$. We will also write $[x_1, x_2, \dots, x_n]^\T$, for the corresponding column vector, where ${}^\T$ denotes the \textit{transpose}.

\begin{xattention}
To avoid confusion, we will use the convention from now on that a generic vector $\v{x}$ is always represented via the standard basis as a column vector. The corresponding row vector is denoted by $\v{x}^\T$.
\end{xattention}

A \textit{matrix} can be viewed as an array of $m$ rows and $n$ columns that defines a \textit{linear transformation} from $\R^n$ to $\R^m$ (or for complex matrices, from $\C^n$ to $\C^m$). The matrix is said to be \textit{square} if $m = n$. If $\v{a}_1, \v{a}_2, \dots, \v{a}_n$ are the columns of $A$, that is, $A = [\v{a}_1, \v{a}_2, \dots, \v{a}_n]$, and if $\v{x} = [x_1, \dots, x_n]^\T$, then $\v{A}\v{x} = x_1 \v{a}_1 + \dots + x_n \v{a}_n$. In particular, the standard basis vector $\v{e}_k$ is mapped to the vector $\v{a}_k, k = 1, \dots, n$. We sometimes use the notation $A = [a_{ij}]$, to denote a matrix whose $(i, j)$-th element is $a_{ij}$. When we wish to emphasize that a matrix $\v{A}$ is real-valued with $m$ rows and $n$ columns, we write $A \in \R^{m \times n}$. The \textit{rank} of a matrix is the number of linearly independent rows or, equivalently, the number of linearly independent columns.

The \textit{transpose} of a matrix $\v{A} = [a_{ij}]$ is the matrix $\v{A}^\T = [a_{ji}]$; that is, the $(i, j)$-th element of $\v{A}^\T$ is the $(j, i)$-th element of $\v{A}$. The \textit{trace} of a square matrix is the sum of its diagonal elements. A useful result is the following cyclic property.

\begin{thm}{Cyclic Property}{CyclicProperty}
The trace is invariant under cyclic permutations: $\tr(\v{ABC}) = \tr(\v{BCA}) = \tr(\v{CAB})$.
\end{thm}

\begin{proof}
It suffices to show that $\tr(\v{DE})$ is equal to $\tr(\v{ED})$ for any $m \times n$ matrix $\v{D} = [d_{ij}]$ and $n \times m$ matrix $\v{E} = [e_{ij}]$. The diagonal elements of $\v{DE}$ are $\sum_{j=1}^n d_{ij} e_{ji}, i = 1, \dots, m$ and the diagonal elements of $\v{ED}$ are $\sum_{i=1}^m e_{ji} d_{ij}, j = 1, \dots, n$. They sum up to the same number $\sum_{i=1}^m \sum_{j=1}^n d_{ij} e_{ji}$.
\end{proof}

A square matrix has an inverse if and only if its columns (or rows) are \textit{linearly independent}. This is the same as the matrix being of \textit{full rank}; that is, its rank is equal to the number of columns. An equivalent statement is that its determinant is not zero. The \textit{determinant} of an $n \times n$ matrix $\v{A} = [a_{i,j}]$ is defined as
\begin{equation}
    \det(\v{A}) \coloneqq \sum_{\vect{\pi}} (-1)^{\zeta(\vect{\pi})} \prod_{i=1}^n a_{\pi_i, i}, \tag{A.3}
\end{equation}
where the sum is over all permutations $\vect{\pi} = (\pi_1, \dots, \pi_n)$ of $(1, \dots, n)$, and $\zeta(\vect{\pi})$ is the number of pairs $(i, j)$ for which $i < j$ and $\pi_i > \pi_j$. For example, $\zeta(2, 3, 4, 1) = 3$ for the pairs $(1, 4), (2, 4), (3, 4)$. The determinant of a \textit{diagonal matrix}, a matrix with only zero elements off the diagonal, is simply the product of its diagonal elements.

Geometrically, the determinant of a square matrix $\v{A} = [\v{a}_1, \dots, \v{a}_n]$ is the (signed) \textit{volume} of the parallelepiped ($n$-dimensional parallelogram) defined by the columns $\v{a}_1, \dots, \v{a}_n$; that is, the set of points $\v{x} = \sum_{i=1}^n \alpha_i \v{a}_i$, where $0 \le \alpha_i \le 1, i = 1, \dots, n$.

The easiest way to compute the determinant of a general matrix is to apply simple operations to the matrix that potentially reduce its complexity (as in the number of non-zero elements, for example), while retaining its determinant:

\begin{itemize}
    \item Adding a multiple of one column (or row) to another does not change the determinant.
    \item Multiplying a column (or row) by a number multiplies the determinant by the same number.
    \item Swapping two rows changes the sign of the determinant.
\end{itemize}

By applying these rules repeatedly, one can reduce any matrix to a diagonal matrix. It follows then that the determinant of the original matrix is equal to the product of the diagonal elements of the resulting diagonal matrix multiplied by a known constant.

\begin{thm}{Useful Matrix Rules}{UsefulMatrixRules}
\begin{enumerate}
    \item $(\v{AB})^\T = \v{B}^\T \v{A}^\T$
    \item $(\v{AB})^{-1} = \v{B}^{-1} \v{A}^{-1}$
    \item $(\v{A}^{-1})^\T = (\v{A}^\T)^{-1} \coloneqq \v{A}^{-\T}$
    \item $\det(\v{AB}) = \det(\v{A}) \det(\v{B})$
    \item $\bx^\T \v{A} \bx = \tr(\v{A} \bx \bx^\T)$
    \item $\det(\v{A}) = \prod_i a_{ii}$ if $\v{A} = [a_{ij}]$ is triangular
\end{enumerate}
\end{thm}

\ref{thm:UsefulMatrixRules} summarises a number of useful matrix rules for the concepts that we have discussed so far.

Next, consider an $n \times p$ matrix $\v{A}$ for which the matrix inverse fails to exist. That is, $\v{A}$ is either non-square ($n \neq p$) or its determinant is 0. Instead of the inverse, we can use its so-called pseudo-inverse, which always exists.

\begin{defn}{Moore--Penrose Pseudo-Inverse}{MoorePenrosePseudoInverse}
The \textit{Moore--Penrose pseudo-inverse} of a real matrix $\v{A} \in \R^{n \times p}$ is defined as the unique matrix $\v{A}^+ \in \R^{p \times n}$ that satisfies the conditions:
\begin{enumerate}
    \item $\v{A} \v{A}^+ \v{A} = \v{A}$
    \item $\v{A}^+ \v{A} \v{A}^+ = \v{A}^+$
    \item $(\v{A} \v{A}^+)^\T = \v{A} \v{A}^+$
    \item $(\v{A}^+ \v{A})^\T = \v{A}^+ \v{A}$
\end{enumerate}
\end{defn}

We can write $\v{A}^+$ explicitly in terms of $\v{A}$ when $\v{A}$ has a full column or row rank. For example, we always have
\begin{equation}
    \v{A}^\T \v{A} \v{A}^+ = \v{A}^\T (\v{A} \v{A}^+)^\T = ((\v{A} \v{A}^+) \v{A})^\T = (\v{A})^\T = \v{A}^\T. \label{eq:PseudoInverseIdentity}
\end{equation}
If $\v{A}$ has a full column rank $p$, then $(\v{A}^\T \v{A})^{-1}$ exists, so that from \eqref{eq:PseudoInverseIdentity} it follows that $\v{A}^+ = (\v{A}^\T \v{A})^{-1} \v{A}^\T$. This is referred to as the \textit{left pseudo-inverse}, as $\v{A}^+ \v{A} = \v{I}_p$. Similarly, if $\v{A}$ has a full row rank $n$, that is, $(\v{A} \v{A}^\T)^{-1}$ exists, then it follows from
\begin{equation*}
    \v{A}^+ \v{A} \v{A}^\T = (\v{A}^+ \v{A})^\T \v{A}^\T = (\v{A} (\v{A}^+ \v{A}))^\T = \v{A}^\T
\end{equation*}
that $\v{A}^+ = \v{A}^\T (\v{A} \v{A}^\T)^{-1}$. This is the \textit{right pseudo-inverse}, as $\v{A} \v{A}^+ = \v{I}_n$. Finally, if $\v{A}$ is of full rank and square, then $\v{A}^+ = \v{A}^{-1}$.

\subsection{Inner Product}
The (Euclidean) \textit{inner product} of two real vectors $\v{x} = [x_1, \dots, x_n]^\T$ and $\v{y} = [y_1, \dots, y_n]^\T$ is defined as the number
\begin{equation*}
    \langle \v{x}, \v{y} \rangle = \sum_{i=1}^n x_i y_i = \v{x}^\T \v{y}.
\end{equation*}
Here $\v{x}^\T \v{y}$ is the matrix multiplication of the $(1 \times n)$ matrix $\v{x}^\T$ and the $(n \times 1)$ matrix $\v{y}$. The inner product induces a geometry on the linear space $\R^n$, allowing for the definition of length, angle, and so on. The inner product satisfies the following properties:
\begin{enumerate}
    \item $\langle \alpha \v{x} + \beta \v{y}, \v{z} \rangle = \alpha \langle \v{x}, \v{z} \rangle + \beta \langle \v{y}, \v{z} \rangle$;
    \item $\langle \v{x}, \v{y} \rangle = \langle \v{y}, \v{x} \rangle$;
    \item $\langle \v{x}, \v{x} \rangle \ge 0$;
    \item $\langle \v{x}, \v{x} \rangle = 0$ if and only if $\v{x} = \v{0}$.
\end{enumerate}
Vectors $\v{x}$ and $\v{y}$ are called \textit{perpendicular} (or \textit{orthogonal}) if $\langle \v{x}, \v{y} \rangle = 0$. The \textit{Euclidean norm} (or length) of a vector $\v{x}$ is defined as
\begin{equation*}
    \|\v{x}\| = \sqrt{x_1^2 + \dots + x_n^2} = \sqrt{\langle \v{x}, \v{x} \rangle}.
\end{equation*}
If $\v{x}$ and $\v{y}$ are perpendicular, then \textit{Pythagoras' theorem} holds:
\begin{equation}
    \|\v{x} + \v{y}\|^2 = \langle \v{x} + \v{y}, \v{x} + \v{y} \rangle = \langle \v{x}, \v{x} \rangle + 2 \langle \v{x}, \v{y} \rangle + \langle \v{y}, \v{y} \rangle = \|\v{x}\|^2 + \|\v{y}\|^2. \tag{A.5}
\end{equation}
A basis $\{\v{\nu}_1, \dots, \v{\nu}_n\}$ of $\R^n$ in which all the vectors are pairwise perpendicular and have norm 1 is called an \textit{orthonormal} (short for orthogonal and normalized) basis. For example, the standard basis is orthonormal.

\begin{thm}{Orthonormal Basis Representation}{OrthonormalBasisRepresentation}
If $\{\v{\nu}_1, \dots, \v{\nu}_n\}$ is an orthonormal basis of $\R^n$, then any vector $\v{x} \in \R^n$ can be expressed as
\begin{equation}
    \v{x} = \langle \v{x}, \v{\nu}_1 \rangle \v{\nu}_1 + \dots + \langle \v{x}, \v{\nu}_n \rangle \v{\nu}_n. \tag{A.6}
\end{equation}
\end{thm}
\begin{proof}
Observe that, because the $\{\v{\nu}_i\}$ form a basis, there exist unique $\alpha_1, \dots, \alpha_n$ such that $\v{x} = \alpha_1 \v{\nu}_1 + \dots + \alpha_n \v{\nu}_n$. By the linearity of the inner product and the orthonormality of the $\{\v{\nu}_i\}$ it follows that $\langle \v{x}, \v{\nu}_j \rangle = \langle \sum_i \alpha_i \v{\nu}_i, \v{\nu}_j \rangle = \alpha_j$.
\end{proof}
An $n \times n$ matrix $\v{V}$ whose columns form an orthonormal basis is called an \textit{orthogonal matrix}.\footnote{The qualifier ``orthogonal'' for such matrices has been fixed by history. A better term would have been ``orthonormal''.} Note that for an orthogonal matrix $\v{V} = [\v{\nu}_1, \dots, \v{\nu}_n]$, we have
\begin{equation*}
    \v{V}^\T \v{V} =
    \begin{bmatrix}
        \v{\nu}_1^\T \\
        \v{\nu}_2^\T \\
        \vdots \\
        \v{\nu}_n^\T
    \end{bmatrix}
    [\v{\nu}_1, \v{\nu}_2, \dots, \v{\nu}_n] =
    \begin{bmatrix}
        \v{\nu}_1^\T \v{\nu}_1 & \v{\nu}_1^\T \v{\nu}_2 & \dots & \v{\nu}_1^\T \v{\nu}_n \\
        \vdots & \vdots & \vdots & \vdots \\
        \v{\nu}_n^\T \v{\nu}_1 & \v{\nu}_n^\T \v{\nu}_2 & \dots & \v{\nu}_n^\T \v{\nu}_n
    \end{bmatrix}
    = \v{I}_n.
\end{equation*}
Hence, $\v{V}^{-1} = \v{V}^\T$. Note also that an orthogonal transformation is \textit{length preserving}; that is, $\v{Vx}$ has the same length as $\v{x}$. This follows from
\begin{equation*}
    \|\v{Vx}\|^2 = \langle \v{Vx}, \v{Vx} \rangle = \v{x}^\T \v{V}^\T \v{V} \v{x} = \v{x}^\T \v{x} = \|\v{x}\|^2.
\end{equation*}
\subsection{Complex Vectors and Matrices}\label{sec:ComplexVectorsMatrices}
Instead of the vector space $\R^n$ of $n$-dimensional real vectors, it is sometimes useful to consider the vector space $\C^n$ of $n$-dimensional complex vectors. In this case the \textit{adjoint} or \textit{conjugate transpose} operation ($*$) replaces the transpose operation ($\T$). This involves the usual transposition of the matrix or vector with the additional step that any complex number $z = x + \ii y$ is replaced by its complex conjugate $\overline{z} = x - \ii y$. For example, if
\begin{equation*}
    \v{x} =
    \begin{bmatrix}
        a_1 + \ii b_1 \\
        a_2 + \ii b_2
    \end{bmatrix}
    \qquad \text{and} \qquad
    \v{A} =
    \begin{bmatrix}
        a_{11} + \ii b_{11} & a_{12} + \ii b_{12} \\
        a_{21} + \ii b_{21} & a_{22} + \ii b_{22}
    \end{bmatrix},
\end{equation*}
then
\begin{equation*}
    \v{x}^* = [a_1 - \ii b_1, a_2 - \ii b_2]
    \qquad \text{and} \qquad
    \v{A}^* =
    \begin{bmatrix}
        a_{11} - \ii b_{11} & a_{21} - \ii b_{21} \\
        a_{12} - \ii b_{12} & a_{22} - \ii b_{22}
    \end{bmatrix}.
\end{equation*}
The (Euclidean) inner product of $\v{x}$ and $\v{y}$ (viewed as column vectors) is now defined as
\begin{equation*}
    \langle \v{x}, \v{y} \rangle = \v{y}^* \v{x} = \sum_{i=1}^n x_i \overline{y}_i,
\end{equation*}
which is no longer symmetric: $\langle \v{x}, \v{y} \rangle = \overline{\langle \v{y}, \v{x} \rangle}$. Note that this generalizes the real-valued inner product. The determinant of a complex matrix $\v{A}$ is defined exactly as in \eqref{eq:DeterminantDef}. As a consequence, $\det(\v{A}^*) = \overline{\det(\v{A})}$.

A complex matrix is said to be \textit{Hermitian} or \textit{self-adjoint} if $\v{A}^* = \v{A}$, and \textit{unitary} if $\v{A}^* \v{A} = \v{I}$ (that is, if $\v{A}^* = \v{A}^{-1}$). For real matrices ``Hermitian'' is the same as ``symmetric'', and ``unitary'' is the same as ``orthogonal''.

\subsection{Orthogonal Projections}
Let $\{\v{u}_1, \dots, \v{u}_k\}$ be a set of linearly independent vectors in $\R^n$. The set
\begin{equation*}
    \mathcal{V} = \Span{\v{u}_1, \dots, \v{u}_k} = \{\alpha_1 \v{u}_1 + \dots + \alpha_k \v{u}_k, \alpha_1, \dots, \alpha_k \in \R\},
\end{equation*}
is called the \textit{linear subspace} spanned by $\{\v{u}_1, \dots, \v{u}_k\}$. The \textit{orthogonal complement} of $\mathcal{V}$, denoted by $\mathcal{V}^\perp$, is the set of all vectors $\v{w}$ that are orthogonal to $\mathcal{V}$, in the sense that $\langle \v{w}, \v{v} \rangle = 0$ for all $\v{v} \in \mathcal{V}$. The matrix $\v{P}$ such that $\v{Px} = \v{x}$, for all $\v{x} \in \mathcal{V}$, and $\v{Px} = \v{0}$, for all $\v{x} \in \mathcal{V}^\perp$ is called the \textit{orthogonal projection matrix} onto $\mathcal{V}$. Suppose that $\v{U} = [\v{u}_1, \dots, \v{u}_k]$ has full rank, in which case $\v{U}^\T \v{U}$ is an invertible matrix. The orthogonal projection matrix $\v{P}$ onto $\mathcal{V} = \Span{\v{u}_1, \dots, \v{u}_k}$ is then given by
\begin{equation*}
    \v{P} = \v{U}(\v{U}^\T \v{U})^{-1}\v{U}^\T.
\end{equation*}
Namely, since $\v{PU} = \v{U}$, the matrix $\v{P}$ projects any vector in $\mathcal{V}$ onto itself. Moreover, $\v{P}$ projects any vector in $\mathcal{V}^\perp$ onto the zero vector. Using the pseudo-inverse, it is possible to specify the projection matrix also for the case where $\v{U}$ not of full rank, leading to the following theorem.

\begin{thm}{Orthogonal Projection}{OrthogonalProjection}
Let $\v{U} = [\v{u}_1, \dots, \v{u}_k]$. Then, the orthogonal projection matrix $\v{P}$ onto $\mathcal{V} = \Span{\v{u}_1, \dots, \v{u}_k}$ is given by
\begin{equation}
    \v{P} = \v{U} \v{U}^+, \tag{A.7}
\end{equation}
where $\v{U}^+$ is the (right) pseudo-inverse of $\v{U}$.
\end{thm}

\begin{proof}
By Property 1 of \ref{defn:MoorePenrosePseudoInverse} we have $\v{PU} = \v{U U}^+ \v{U} = \v{U}$, so that $\v{P}$ projects any vector in $\mathcal{V}$ onto itself. Moreover, $\v{P}$ projects any vector in $\mathcal{V}^\perp$ onto the zero vector.
\end{proof}

Note that in the special case where $\v{u}_1, \dots, \v{u}_k$ above form an orthonormal basis of $\mathcal{V}$, then the projection onto $\mathcal{V}$ is very simple to describe, namely we have
\begin{equation}
    \v{Px} = \v{U U}^\T \v{x} = \sum_{i=1}^k \langle \v{x}, \v{u}_i \rangle \v{u}_i. \tag{A.8}
    \label{eq:OrthonormalProjectionSum}
\end{equation}

For any point $\v{x} \in \R^n$, the point in $\mathcal{V}$ that is closest to $\v{x}$ is its orthogonal projection $\v{Px}$, as the following theorem shows.

\begin{thm}{Orthogonal Projection and Minimal Distance}{OrthogonalProjectionAndMinimalDistance}
Let $\{\v{u}_1, \dots, \v{u}_k\}$ be an orthonormal basis of subspace $\mathcal{V}$ and let $\v{P}$ be the orthogonal projection matrix onto $\mathcal{V}$. The solution to the minimization program
\begin{equation*}
    \min_{\v{y} \in \mathcal{V}} \|\v{x} - \v{y}\|^2
\end{equation*}
is $\v{y} = \v{Px}$. That is, $\v{Px} \in \mathcal{V}$ is closest to $\v{x}$.
\end{thm}

\begin{proof}
We can write each point $\v{y} \in \mathcal{V}$ as $\v{y} = \sum_{i=1}^k \alpha_i \v{u}_i$. Consequently,
\begin{equation*}
    \|\v{y} - \v{x}\|^2 = \left\langle \v{x} - \sum_{i=1}^k \alpha_i \v{u}_i, \v{x} - \sum_{i=1}^k \alpha_i \v{u}_i \right\rangle = \|\v{x}\|^2 - 2 \sum_{i=1}^k \alpha_i \langle \v{x}, \v{u}_i \rangle + \sum_{i=1}^k \alpha_i^2.
\end{equation*}
Minimizing this with respect to the $\{\alpha_i\}$ gives $\alpha_i = \langle \v{x}, \v{u}_i \rangle, i = 1, \dots, k$. In view of \cref{eq:OrthonormalProjectionSum}, the optimal $\v{y}$ is thus $\v{Px}$.
\end{proof}

\subsection{Eigenvalues and Eigenvectors}
Let $\v{A}$ be an $n \times n$ matrix. If $\v{A}\v{v} = \lambda \v{v}$ for some number $\lambda$ and non-zero vector $\v{v}$, then $\lambda$ is called an \textit{eigenvalue} of $\v{A}$ with \textit{eigenvector} $\v{v}$.

If $(\lambda, \v{v})$ is an (eigenvalue, eigenvector) pair, the matrix $\lambda \v{I} - \v{A}$ maps any multiple of $\v{v}$ to the zero vector. Consequently, the columns of $\lambda \v{I} - \v{A}$ are linearly dependent, and hence their determinant is 0. This provides a way to identify the eigenvalues, namely as the $r \le n$ different roots $\lambda_1, \dots, \lambda_r$ of the \textit{characteristic polynomial}
\begin{equation*}
    \det(\lambda \v{I} - \v{A}) = (\lambda - \lambda_1)^{\alpha_1} \cdots (\lambda - \lambda_r)^{\alpha_r},
\end{equation*}
where $\alpha_1 + \dots + \alpha_r = n$. The integer $\alpha_i$ is called the \textit{algebraic multiplicity} of $\lambda_i$. The eigenvectors that correspond to an eigenvalue $\lambda_i$ lie in the \textit{kernel} or \textit{null space} of the matrix $\lambda_i \v{I} - \v{A}$; that is, the linear space of vectors $\v{v}$ such that $(\lambda_i \v{I} - \v{A})\v{v} = \v{0}$. This space is called the \textit{eigenspace} of $\lambda_i$. Its dimension, $d_i \in \{1, \dots, n\}$, is called the \textit{geometric multiplicity} of $\lambda_i$. It always holds that $d_i \le \alpha_i$. If $\sum_i d_i = n$, then we can construct a basis for $\R^n$ consisting of eigenvectors.

A matrix for which the algebraic and geometric multiplicities of all its eigenvalues are the same is called \textit{semi-simple}. This is equivalent to the matrix being \textit{diagonalizable}, meaning that there is a matrix $\v{V}$ and a diagonal matrix $\v{D}$ such that
\begin{equation*}
    \v{A} = \v{V} \v{D} \v{V}^{-1}.
\end{equation*}
To see that this so-called \textit{eigen-decomposition} holds, suppose $\v{A}$ is a semi-simple matrix with eigenvalues
\begin{equation*}
    \underbrace{\lambda_1, \dots, \lambda_1}_{d_1}, \dots, \underbrace{\lambda_r, \dots, \lambda_r}_{d_r}.
\end{equation*}
Let $\v{D}$ be the diagonal matrix whose diagonal elements are the eigenvalues of $\v{A}$, and let $\v{V}$ be a matrix whose columns are linearly independent eigenvectors corresponding to these eigenvalues. Then, for each (eigenvalue, eigenvector) pair $(\lambda, \v{v})$, we have $\v{A}\v{v} = \lambda \v{v}$. Hence, in matrix notation, we have $\v{A}\v{V} = \v{V}\v{D}$, and so $\v{A} = \v{V}\v{D}\v{V}^{-1}$.

\subsection{Matrix Decompositions}
Matrix decompositions are frequently employed in linear algebra to simplify proofs, mitigate numerical instability, and expedite computations. We mention three important matrix decompositions: (P)LU, QR, and SVD.

\subsubsection{(P)LU Decomposition}
Every invertible matrix $\v{A}$ can be written as the product of three matrices:
\begin{equation}
    \v{A} = \v{PLU}, \label{eq:PLUDecomposition}
\end{equation}
where $\v{L}$ is a lower triangular matrix, $\v{U}$ an upper triangular matrix, and $\v{P}$ a \textit{permutation matrix}. A permutation matrix is a square matrix with a single 1 in each row and column, and zeros otherwise. The matrix product $\v{PB}$ simply permutes the rows of a matrix $\v{B}$ and, likewise, $\v{BP}$ permutes its columns. A decomposition of the form \cref{eq:PLUDecomposition} is called a \textit{PLU decomposition}. As a permutation matrix is orthogonal, its transpose is equal to its inverse, and so we can write \cref{eq:PLUDecomposition} as
\begin{equation*}
    \v{P}^\T \v{A} = \v{LU}.
\end{equation*}
The decomposition is not unique, and in many cases $\v{P}$ can be taken to be the identity matrix, in which case we speak of the \textit{LU decomposition} of $\v{A}$, also called the LR for left--right (triangular) decomposition.

A PLU decomposition of an invertible $n \times n$ matrix $\v{A}_0$ can be obtained recursively as follows. The first step is to swap the rows of $\v{A}_0$ such that the element in the first column and first row of the pivoted matrix is as large as possible in absolute value. Write the resulting matrix as
\begin{equation*}
    \widetilde{\v{P}}_0 \v{A}_0 =
    \begin{bmatrix}
        a_1 & \v{b}_1^\T \\
        \v{c}_1 & \v{D}_1
    \end{bmatrix},
\end{equation*}
where $\widetilde{\v{P}}_0$ is the permutation matrix that swaps the first and $k$-th row, where $k$ is the row that contains the largest element in the first column. Next, add the matrix $-\v{c}_1 [1, \v{b}_1^\T / a_1]$ to the last $n-1$ rows of $\widetilde{\v{P}}_0 \v{A}_0$, to obtain the matrix
\begin{equation*}
    \begin{bmatrix}
        a_1 & \v{b}_1^\T \\
        \v{0} & \v{D}_1 - \v{c}_1 \v{b}_1^\T / a_1
    \end{bmatrix}
    \coloneqq
    \begin{bmatrix}
        a_1 & \v{b}_1^\T \\
        \v{0} & \v{A}_1
    \end{bmatrix}.
\end{equation*}
In effect, we add some multiple of the first row to each of the remaining rows in order to obtain zeros in the first column, except for the first element.

We now apply the same procedure to $\v{A}_1$ as we did to $\v{A}_0$ and then to subsequent smaller matrices $\v{A}_2, \dots, \v{A}_{n-1}$:
\begin{enumerate}
    \item Swap the first row with the row having the maximal absolute value element in the first column.
    \item Make every other element in the first column equal to 0 by adding appropriate multiples of the first row to the other rows.
\end{enumerate}
Suppose that $\v{A}_t$ has a PLU decomposition $\v{P}_t \v{L}_t \v{U}_t$. Then it is easy to check that
\begin{equation}
    \underbrace{\widetilde{\v{P}}_{t-1}^\T
    \begin{bmatrix}
        1 & \v{0}^\T \\
        \v{0} & \v{P}_t
    \end{bmatrix}}_{\v{P}_{t-1}}
    \underbrace{
    \begin{bmatrix}
        1 & \v{0}^\T \\
        \v{P}_t^\T \v{c}_t / a_t & \v{L}_t
    \end{bmatrix}}_{\v{L}_{t-1}}
    \underbrace{
    \begin{bmatrix}
        a_t & \v{b}_t^\T \\
        \v{0} & \v{U}_t
    \end{bmatrix}}_{\v{U}_{t-1}} \tag{A.10}
\end{equation}
is a PLU decomposition of $\v{A}_{t-1}$. Since the PLU decomposition for the scalar $\v{A}_{n-1}$ is trivial, by working backwards, we obtain a PLU decomposition $\v{P}_0 \v{L}_0 \v{U}_0$ of $\v{A}$.

\subsubsection{Woodburry Identity}
LU (or more generally PLU) decompositions can also be applied to block matrices. A starting point is the following LU decomposition for a general $2 \times 2$ matrix:
\begin{equation*}
    \begin{bmatrix}
        a & b \\
        c & d
    \end{bmatrix}
    =
    \begin{bmatrix}
        a & 0 \\
        c & d - bc/a
    \end{bmatrix}
    \begin{bmatrix}
        1 & b/a \\
        0 & 1
    \end{bmatrix},
\end{equation*}
which holds as long as $a \ne 0$; this can be seen by simply writing out the matrix product. The block matrix generalization for matrices $\v{A} \in \R^{n \times n}, \v{B} \in \R^{n \times k}, \v{C} \in \R^{k \times n}, \v{D} \in \R^{k \times k}$ is
\begin{equation} \label{eq:BlockLU}
    \Sigma \coloneqq \begin{bmatrix}
        \v{A} & \v{B} \\
        \v{C} & \v{D}
    \end{bmatrix}
    =
    \begin{bmatrix}
        \v{A} & \v{O}_{n \times k} \\
        \v{C} & \v{D} - \v{C} \v{A}^{-1} \v{B}
    \end{bmatrix}
    \begin{bmatrix}
        \v{I}_n & \v{A}^{-1} \v{B} \\
        \v{O}_{k \times n} & \v{I}_k
    \end{bmatrix},
\end{equation}
provided that $\v{A}$ is invertible (again, write out the block matrix product). Here, we use the notation $\v{O}_{p \times q}$ to denote the $p \times q$ matrix of zeros. We can further rewrite this as:
\begin{equation*}
    \Sigma =
    \begin{bmatrix}
        \v{I}_n & \v{O}_{n \times k} \\
        \v{C} \v{A}^{-1} & \v{I}_k
    \end{bmatrix}
    \begin{bmatrix}
        \v{A} & \v{O}_{n \times k} \\
        \v{O}_{k \times n} & \v{D} - \v{C} \v{A}^{-1} \v{B}
    \end{bmatrix}
    \begin{bmatrix}
        \v{I}_n & \v{A}^{-1} \v{B} \\
        \v{O}_{k \times n} & \v{I}_k
    \end{bmatrix}.
\end{equation*}
Thus, inverting both sides, we obtain
\begin{equation*}
    \Sigma^{-1} =
    \begin{bmatrix}
        \v{I}_n & \v{A}^{-1} \v{B} \\
        \v{O}_{k \times n} & \v{I}_k
    \end{bmatrix}^{-1}
    \begin{bmatrix}
        \v{A} & \v{O}_{n \times k} \\
        \v{O}_{k \times n} & \v{D} - \v{C} \v{A}^{-1} \v{B}
    \end{bmatrix}^{-1}
    \begin{bmatrix}
        \v{I}_n & \v{O}_{n \times k} \\
        \v{C} \v{A}^{-1} & \v{I}_k
    \end{bmatrix}^{-1}.
\end{equation*}
Inversion of the above block matrices gives (again write out)
\begin{equation} \label{eq:BlockInverseIntermed}
    \Sigma^{-1} =
    \begin{bmatrix}
        \v{I}_n & -\v{A}^{-1} \v{B} \\
        \v{O}_{k \times n} & \v{I}_k
    \end{bmatrix}
    \begin{bmatrix}
        \v{A}^{-1} & \v{O}_{n \times k} \\
        \v{O}_{k \times n} & (\v{D} - \v{C} \v{A}^{-1} \v{B})^{-1}
    \end{bmatrix}
    \begin{bmatrix}
        \v{I}_n & \v{O}_{n \times k} \\
        -\v{C} \v{A}^{-1} & \v{I}_k
    \end{bmatrix}.
\end{equation}
Assuming that $\v{D}$ is invertible, we could also perform a block UL (as opposed to LU) decomposition:
\begin{equation} \label{eq:BlockUL}
    \Sigma =
    \begin{bmatrix}
        \v{A} - \v{B} \v{D}^{-1} \v{C} & \v{B} \\
        \v{O}_{k \times n} & \v{D}
    \end{bmatrix}
    \begin{bmatrix}
        \v{I}_n & \v{O}_{n \times k} \\
        \v{D}^{-1} \v{C} & \v{I}_k
    \end{bmatrix},
\end{equation}
which, after a similar calculation as the one above, yields
\begin{equation} \label{eq:BlockInverseUL}
    \Sigma^{-1} =
    \begin{bmatrix}
        \v{I}_n & \v{O}_{n \times k} \\
        -\v{D}^{-1} \v{C} & \v{I}_k
    \end{bmatrix}
    \begin{bmatrix}
        (\v{A} - \v{B} \v{D}^{-1} \v{C})^{-1} & \v{O}_{n \times k} \\
        \v{O}_{k \times n} & \v{D}^{-1}
    \end{bmatrix}
    \begin{bmatrix}
        \v{I}_n & -\v{B} \v{D}^{-1} \\
        \v{O}_{k \times n} & \v{I}_k
    \end{bmatrix}.
\end{equation}
The upper-left block of $\Sigma^{-1}$ from \cref{eq:BlockInverseUL} must be the same as the upper-left block of $\Sigma^{-1}$ from \cref{eq:BlockInverseIntermed}, leading to the \textit{Woodbury identity}:
\begin{equation} \label{eq:WoodburyIdentity}
    (\v{A} - \v{B} \v{D}^{-1} \v{C})^{-1} = \v{A}^{-1} + \v{A}^{-1} \v{B} (\v{D} - \v{C} \v{A}^{-1} \v{B})^{-1} \v{C} \v{A}^{-1}.
\end{equation}
From \cref{eq:BlockLU} and the fact that the determinant of a product is the product of the determinants, we see that $\det(\Sigma) = \det(\v{A}) \det(\v{D} - \v{C} \v{A}^{-1} \v{B})$. Similarly, from \cref{eq:BlockUL} we have $\det(\Sigma) = \det(\v{A} - \v{B} \v{D}^{-1} \v{C}) \det(\v{D})$, leading to the identity
\begin{equation} \label{eq:DeterminantIdentity}
    \det(\v{A} - \v{B} \v{D}^{-1} \v{C}) \det(\v{D}) = \det(\v{A}) \det(\v{D} - \v{C} \v{A}^{-1} \v{B}).
\end{equation}

\subsubsection{Cholesky Decomposition}

If $\bA$ is a real-valued positive definite matrix (and therefore symmetric), e.g., a covariance matrix, then an LU decomposition can be achieved with matrices $\bL$ and $\bU = \bL^\T$.

\begin{thm}{Cholesky Decomposition}{CholeskyDecomposition}
A real-valued positive definite matrix $\bA = [a_{ij}] \in \R^{n \times n}$ can be decomposed as
\begin{equation}
    \label{CholeskyFactorization}
    \bA = \bL \bL^\T,
\end{equation}
where the real $n \times n$ lower triangular matrix $\bL = [l_{kj}]$ satisfies the recursive formula
\begin{equation}
    \label{CholeskyRecursiveFormula}
    l_{kj} = \frac{a_{kj} - \sum_{i=1}^{j-1} l_{ji} l_{ki}}{\sqrt{a_{jj} - \sum_{i=1}^{j-1} l_{ji}^2}}, \quad \text{where} \quad \sum_{i=1}^0 l_{ji} l_{ki} := 0
\end{equation}
for $k = 1, \dots, n$ and $j = 1, \dots, k$.
\end{thm}

\begin{proof}
The proof is by inductive construction. For $k = 1, \dots, n$, let $\bA_k$ be the left-upper $k \times k$ submatrix of $\bA = \bA_n$. With $\be_1 := [1, 0, \dots, 0]^\T$, we have $\bA_1 = a_{11} = \be_1^\T \bA \be_1 > 0$ by the positive-definiteness of $\bA$. It follows that $l_{11} = \sqrt{a_{11}}$. Suppose that $\bA_{k-1}$ has a Cholesky factorization $\bL_{k-1}\bL_{k-1}^\T$ with $\bL_{k-1}$ having strictly positive diagonal elements, we can construct a Cholesky factorization of $\bA_k$ as follows. First write
\begin{equation}
    \label{MatrixAkBlock}
    \bA_k = \begin{bmatrix} \bL_{k-1}\bL_{k-1}^\T & \ba_{k-1} \\ \ba_{k-1}^\T & a_{kk} \end{bmatrix}
\end{equation}
and propose $\bL_k$ to be of the form
\begin{equation}
    \label{MatrixLkBlock}
    \bL_k = \begin{bmatrix} \bL_{k-1} & \v{0} \\ \v{l}_k^\T & l_{kk} \end{bmatrix}
\end{equation}
for some vector $\v{l}_k \in \R^{k-1}$ and scalar $l_{kk}$, for which it must hold that
\begin{equation}
    \label{MatrixBlockMultiplication}
    \begin{bmatrix} \bL_{k-1}\bL_{k-1}^\T & \ba_{k-1} \\ \ba_{k-1}^\T & a_{kk} \end{bmatrix} = \begin{bmatrix} \bL_{k-1} & \v{0} \\ \v{l}_k^\T & l_{kk} \end{bmatrix} \begin{bmatrix} \bL_{k-1}^\T & \v{l}_k \\ \v{0}^\T & l_{kk} \end{bmatrix}.
\end{equation}
To establish that such an $\v{l}_k$ and $l_{kk}$ exist, we must verify that the set of equations
\begin{equation}
    \label{CholeskyEquationSystem}
    \begin{aligned}
        \bL_{k-1} \v{l}_k &= \ba_{k-1} \\
        \v{l}_k^\T \v{l}_k + l_{kk}^2 &= a_{kk}
    \end{aligned}
\end{equation}
has a solution. The system $\bL_{k-1} \v{l}_k = \ba_{k-1}$ has a unique solution, because (by assumption) $\bL_{k-1}$ is lower diagonal with strictly positive entries down the main diagonal and we can solve for $\v{l}_k$ using forward substitution: $\v{l}_k = \bL_{k-1}^{-1} \ba_{k-1}$. We can solve the second equation as $l_{kk} = \sqrt{a_{kk} - \|\v{l}_k\|^2}$, provided that the term within the square root is positive. We demonstrate this using the fact that $\bA$ is a positive definite matrix. In particular, for $\bx \in \R^n$ of the form $[\bx_1^\T, x_2, \v{0}^\T]^\T$, where $\bx_1$ is a non-zero $(k-1)$-dimensional vector and $x_2$ a non-zero number, we have
\begin{equation}
    \label{PositiveDefiniteCheck}
    0 < \bx^\T \bA \bx = [\bx_1^\T, x_2] \begin{bmatrix} \bL_{k-1}\bL_{k-1}^\T & \ba_{k-1} \\ \ba_{k-1}^\T & a_{kk} \end{bmatrix} \begin{bmatrix} \bx_1 \\ x_2 \end{bmatrix} = \|\bL_{k-1}^\T \bx_1\|^2 + 2\bx_1^\T \ba_{k-1} x_2 + a_{kk} x_2^2.
\end{equation}
Now take $\bx_1 = -x_2 \bL_{k-1}^{-\T} \v{l}_k$ to obtain $0 < \bx^\T \bA \bx = x_2^2 (a_{kk} - \|\v{l}_k\|^2)$. Therefore, (\ref{CholeskyEquationSystem}) can be uniquely solved. As we have already solved it for $k=1$, we can solve it for any $k=1, \dots, n$, leading to the recursive formula (\ref{CholeskyRecursiveFormula}) and Algorithm \ref{AlgCholeskyDecomposition} below.
\end{proof}

An implementation of Cholesky's decomposition, using the notation from the proof of \ref{thm:CholeskyDecomposition}, is as follows. The running cost of this algorithm is $\bigO(n^3)$.

\begin{algorithm}[H]
\caption{Cholesky Decomposition}
\label{AlgCholeskyDecomposition}
\begin{algorithmic}[1]
    \Require Positive-definite $n \times n$ matrix $\bA_n$ with entries $\{a_{ij}\}$.
    \Ensure Lower triangular $\bL_n$ such that $\bL_n\bL_n^\T = \bA_n$.
    \State $\bL_1 \la \sqrt{a_{11}}$
    \For{$k = 2, \dots, n$}
        \State $\v{a}_{k-1} \la [a_{1k}, \dots, a_{k-1,k}]^\T$
        \State $\v{l}_{k-1} \la \bL_{k-1}^{-1} \v{a}_{k-1}$
        \State $l_{kk} \la \sqrt{a_{kk} - \v{l}_{k-1}^\T \v{l}_{k-1}}$
        \State $\bL_k \la \begin{bmatrix} \bL_{k-1} & \v{0} \\ \v{l}_{k-1}^\T & l_{kk} \end{bmatrix}$
    \EndFor
    \State \Return $\bL_n$
\end{algorithmic}
\end{algorithm}

\subsection{Functional Analysis}
Much of the previous theory on Euclidean vector spaces can be generalized to vector spaces of \textit{functions}. Every element of a (real-valued) \textit{function space} $\scH$ is a function from some set $\mathcal{X}$ to $\R$, and elements can be added and scalar multiplied as if they were vectors. In other words, if $f \in \scH$ and $g \in \scH$, then $\alpha f + \beta g \in \scH$ for all $\alpha, \beta \in \R$. On $\scH$ we can impose an inner product as a mapping $\langle \cdot, \cdot \rangle$ from $\scH \times \scH$ to $\R$ that satisfies
\begin{enumerate}
    \item $\langle \alpha f_1 + \beta f_2, g \rangle = \alpha \langle f_1, g \rangle + \beta \langle f_2, g \rangle$;
    \item $\langle f, g \rangle = \langle g, f \rangle$;
    \item $\langle f, f \rangle \ge 0$;
    \item $\langle f, f \rangle = 0$ if and only if $f = 0$ (the zero function).
\end{enumerate}
We focus on real-valued function spaces, although the theory for complex-valued function spaces is similar (and sometimes easier), under suitable modifications (e.g., $\langle f, g \rangle = \overline{\langle g, f \rangle}$).

We say that two elements $f$ and $g$ in $\scH$ are \textit{orthogonal} to each other with respect to this inner product if $\langle f, g \rangle = 0$. Given an inner product, we can measure distances between elements of the function space $\scH$ using the \textit{norm}
\begin{equation*}
    \|f\| \coloneqq \sqrt{\langle f, f \rangle}.
\end{equation*}
For example, the distance between two functions $f_m$ and $f_n$ is given by $\|f_m - f_n\|$. The space $\scH$ is said to be \textit{complete} if every sequence of functions $f_1, f_2, \dots \in \scH$ for which
\begin{equation}
    \|f_m - f_n\| \to 0 \text{ as } m, n \to \infty, \label{eq:CauchySequence}
\end{equation}
converges to some $f \in \scH$; that is, $\|f - f_n\| \to 0$ as $n \to \infty$. A sequence that satisfies \eqref{eq:CauchySequence} is called a \textit{Cauchy sequence}.

A complete inner product space is called a \textit{Hilbert space}. The most fundamental Hilbert space of functions is the space $L^2$. An in-depth introduction to $L^2$ requires some measure theory [6]. For our purposes, it suffices to assume that $\mathcal{X} \subseteq \R^d$ and that on $\mathcal{X}$ a \textit{measure} $\mu$ is defined which assigns to each suitable\footnote{Not all sets have a measure. Suitable sets are Borel sets, which can be thought of as countable unions of rectangles.} set $A$ a positive number $\mu(A) \ge 0$ (e.g., its volume). In many cases of interest $\mu$ is of the form
\begin{equation}
    \mu(A) = \int_A w(x) \di x \label{eq:MeasureDensityForm}
\end{equation}
where $w \ge 0$ is a positive function on $\mathcal{X}$ which is called the \textit{density} of $\mu$ with respect to the Lebesgue measure (the natural volume measure on $\R^d$). We write $\mu(\di x) = w(x) \di x$ to indicate that $\mu$ has density $w$. Another important case is where
\begin{equation}
    \mu(A) = \sum_{x \in A \cap \Z^d} w(x), \label{eq:DiscreteMeasureForm}
\end{equation}
where $w \ge 0$ is again called the density of $\mu$, but now with respect to the counting measure on $\R^d$ (which counts the points of $\Z^d$). Integrals with respect to measures $\mu$ in \eqref{eq:MeasureDensityForm} and \eqref{eq:DiscreteMeasureForm} can now be defined as
\begin{equation*}
    \int f(x) \mu(\di x) = \int f(x) w(x) \di x,
\end{equation*}
and
\begin{equation*}
    \int f(x) \mu(\di x) = \sum_x f(x) w(x),
\end{equation*}
respectively. We assume for simplicity that $\mu$ has the form \eqref{eq:MeasureDensityForm}. For measures of the form \eqref{eq:DiscreteMeasureForm} (so-called discrete measures), replace integrals by sums in what follows.

\begin{defn}{$L^2$ Space}{L2Space}
Let $\mathcal{X}$ be a subset of $\R^d$ with measure $\mu(\di x) = w(x) \di x$. The Hilbert space $L^2(\mathcal{X}, \mu)$ is the linear space of functions from $\mathcal{X}$ to $\R$ that satisfy
\begin{equation}
    \int_{\mathcal{X}} f(x)^2 w(x) \di x < \infty, \label{eq:L2Integrability}
\end{equation}
and with inner product
\begin{equation}
    \langle f, g \rangle = \int_{\mathcal{X}} f(x) g(x) w(x) \di x. \label{eq:L2InnerProduct}
\end{equation}
\end{defn}

Let $\scH$ be a Hilbert space. A set of functions $\{f_i, i \in \mathcal{I}\}$ is called an \textit{orthonormal system} if
\begin{enumerate}
    \item the norm of every $f_i$ is 1; that is, $\langle f_i, f_i \rangle = 1$ for all $i \in \mathcal{I}$,
    \item the $\{f_i\}$ are orthogonal; that is, $\langle f_i, f_j \rangle = 0$ for $i \neq j$.
\end{enumerate}
It follows then that the $\{f_i\}$ are linearly independent; that is, the only linear combination $\sum_j \alpha_j f_j(x)$ that is equal to $f_i(x)$ for all $x$ is the one where $\alpha_i = 1$ and $\alpha_j = 0$ for $j \neq i$. An orthonormal system $\{f_i\}$ is called an \textit{orthonormal basis} if there is no other $f \in \scH$ that is orthogonal to all the $\{f_i, i \in \mathcal{I}\}$ (other than the zero function). Although the general theory allows for uncountable bases, in practice the set $\mathcal{I}$ is taken to be countable.

We conclude with a number of key results in functional analysis. The first one is the celebrated \textit{Cauchy--Schwarz inequality}.

\begin{thm}{Cauchy--Schwarz}{CauchySchwarz}
Let $\scH$ be a Hilbert space. For every $f, g \in \scH$ it holds that
\begin{equation*}
    |\langle f, g \rangle| \le \|f\| \|g\|.
\end{equation*}
\end{thm}

\begin{proof}
The inequality is trivially true for $g = 0$ (zero function). For $g \neq 0$, we can write $f = \alpha g + h$, where $h \perp g$ and $\alpha = \langle f, g \rangle / \|g\|^2$. Consequently, $\|f\|^2 = |\alpha|^2 \|g\|^2 + \|h\|^2 \ge |\alpha|^2 \|g\|^2$. The result follows after rearranging this last inequality.
\end{proof}

Let $\mathcal{V}$ and $\mathcal{W}$ be two linear vector spaces (for example, Hilbert spaces) on which norms $\|\cdot\|_{\mathcal{V}}$ and $\|\cdot\|_{\mathcal{W}}$ are defined. Suppose $A : \mathcal{V} \to \mathcal{W}$ is a mapping from $\mathcal{V}$ to $\mathcal{W}$. When $\mathcal{W} = \mathcal{V}$, such a mapping is often called an \textit{operator}; when $\mathcal{W} = \R$ it is called a \textit{functional}. Mapping $A$ is said to be \textit{linear} if $A(\alpha f + \beta g) = \alpha A(f) + \beta A(g)$. In this case we write $Af$ instead of $A(f)$. If there exists $\gamma < \infty$ such that
\begin{equation}
    \|Af\|_{\mathcal{W}} \le \gamma \|f\|_{\mathcal{V}}, \quad f \in \mathcal{V}, \label{eq:BoundedOperator}
\end{equation}
then $A$ is said to be a \textit{bounded mapping}. The smallest $\gamma$ for which \eqref{eq:BoundedOperator} holds is called the \textit{norm} of $A$; denoted by $\|A\|$. A (not necessarily linear) mapping $A : \mathcal{V} \to \mathcal{W}$ is said to be \textit{continuous} at $f$ if for any sequence $f_1, f_2, \dots$ converging to $f$ the sequence $A(f_1), A(f_2), \dots$ converges to $A(f)$. That is, if
\begin{equation}
    \forall \varepsilon > 0, \exists \delta > 0 : \forall g \in \mathcal{V}, \|f - g\|_{\mathcal{V}} < \delta \implies \|A(f) - A(g)\|_{\mathcal{W}} < \varepsilon. \label{eq:ContinuousMapping}
\end{equation}
If the above property holds for every $f \in \mathcal{V}$, then the mapping $A$ itself is called \textit{continuous}.

\begin{thm}{Continuity and Boundedness for Linear Mappings}{ContinuityBoundedness}
For a linear mapping, continuity and boundedness are equivalent.
\end{thm}

\begin{proof}
Let $A$ be linear and bounded. We may assume that $A$ is non-zero (otherwise the statement holds trivially), and that therefore $0 < \|A\| < \infty$. Taking $\delta < \varepsilon / \|A\|$ in \eqref{eq:ContinuousMapping} now ensures that $\|Af - Ag\|_{\mathcal{W}} \le \|A\| \|f - g\|_{\mathcal{V}} < \|A\| \delta < \varepsilon$. This shows that $A$ is continuous.
Conversely, suppose $A$ is continuous. In particular, it is continuous at $f = 0$ (the zero element of $\mathcal{V}$). Thus, take $f = 0$ and let $\varepsilon$ and $\delta$ be as in \eqref{eq:ContinuousMapping}. For any $g \neq 0$, let $h = \delta / (2\|g\|_{\mathcal{V}}) g$. As $\|h\|_{\mathcal{V}} = \delta/2 < \delta$, it follows from \eqref{eq:ContinuousMapping} that
\begin{equation*}
    \|Ah\|_{\mathcal{W}} = \frac{\delta}{2\|g\|_{\mathcal{V}}} \|Ag\|_{\mathcal{W}} < \varepsilon.
\end{equation*}
Rearranging the last inequality gives $\|Ag\|_{\mathcal{W}} < 2\varepsilon / \delta \|g\|_{\mathcal{V}}$, showing that $A$ is bounded.
\end{proof}

\begin{thm}{Riesz Representation Theorem}{RieszRepresentation}
Any bounded linear functional $\phi$ on a Hilbert space $\scH$ can be represented as $\phi(h) = \langle h, g \rangle$, for some $g \in \scH$ (depending on $\phi$).
\end{thm}

\begin{proof}
Let $P$ be the projection of $\scH$ onto the nullspace $\mathcal{N}$ of $\phi$; that is, $\mathcal{N} = \{g \in \scH : \phi(g) = 0\}$. If $\phi$ is not the 0-functional, then there exists a $g_0 \neq 0$ with $\phi(g_0) \neq 0$. Let $g_1 = g_0 - Pg_0$. Then $g_1 \perp \mathcal{N}$ and $\phi(g_1) = \phi(g_0)$. Take $g_2 = g_1 / \phi(g_1)$. For any $h \in \scH$, $f \coloneqq h - \phi(h)g_2$ lies in $\mathcal{N}$. As $g_2 \perp \mathcal{N}$ it holds that $\langle f, g_2 \rangle = 0$, which is equivalent to $\langle h, g_2 \rangle = \phi(h) \|g_2\|^2$. By defining $g = g_2 / \|g_2\|^2$ we have found our representation.
\end{proof}

\subsection{Fourier Transforms}
A central concern in mathematical physics and engineering mathematics is the transformation of equations into a coordinate system where expressions simplify, decouple, and become amenable to computation and analysis.

We will now briefly introduce the Fourier transform. Before doing so, we will extend the concept of $L^2$ space of real-valued functions as follows.

\begin{defn}{$L^p$ Space}{LpSpace}
Let $\mathcal{X}$ be a subset of $\R^d$ with measure $\mu(\di x) = w(x) \di x$ and $p \in [1, \infty)$. Then $L^p(\mathcal{X}, \mu)$ is the linear space of functions from $\mathcal{X}$ to $\C$ that satisfy
\begin{equation}
    \int_{\mathcal{X}} |f(x)|^p w(x) \di x < \infty. \label{eq:LpSpaceCondition}
\end{equation}
When $p=2$, $L^2(\mathcal{X}, \mu)$ is in fact a Hilbert space equipped with inner product
\begin{equation}
    \langle f, g \rangle = \int_{\mathcal{X}} f(x) \overline{g(x)} w(x) \di x. \label{eq:LpInnerProduct}
\end{equation}
\end{defn}

\begin{defn}{(Multivariate) Fourier Transform}{MultivariateFourierTransform}
The \textit{Fourier transform} $\mathcal{F}[f]$ of a (real- or complex-valued) function $f \in L^1(\R^d)$ is the function $\widetilde{f}$ defined as
\begin{equation*}
    \widetilde{f}(t) \coloneqq \int_{\R^d} \e^{-\ii 2\pi t^\T x} f(x) \di x, \quad t \in \R^d.
\end{equation*}
\end{defn}

The Fourier transform $\widetilde{f}$ is continuous, uniformly bounded (since $f \in L^1(\R^d)$ implies that $|\widetilde{f}(t)| \le \int_{\R^d} |f(x)| \di x < \infty$), and satisfies $\lim_{\|t\| \to \infty} \widetilde{f}(t) = 0$ (a result known as the \textit{Riemann--Lebesgue lemma}). However, $|\widetilde{f}|$ does not necessarily have a finite integral. A simple example in $\R^1$ is the Fourier transform of $f(x) = \I_{\{-1/2 < x < 1/2\}}$. Then $\widetilde{f}(t) = \sin(\pi t)/(\pi t) = \mathrm{sinc}(\pi t)$, which is not absolutely integrable.

\begin{defn}{(Multivariate) Inverse Fourier Transform}{MultivariateInverseFourierTransform}
The \textit{inverse Fourier transform} $\mathcal{F}^{-1}[\widetilde{f}]$ of a (real- or complex-valued) function $\widetilde{f} \in L^1(\R^d)$ is the function $\breve{f}$ defined as
\begin{equation*}
    \breve{f}(x) \coloneqq \int_{\R^d} \e^{\ii 2\pi t^\T x} \widetilde{f}(t) \di t, \quad x \in \R^d.
\end{equation*}
\end{defn}

As one would hope, it holds that if $f$ and $\mathcal{F}[f]$ are both in $L^1(\R^d)$, then $f = \mathcal{F}^{-1}[\mathcal{F}[f]]$ almost everywhere.
The Fourier transform enjoys many interesting and useful properties, some of which we list below.

\begin{enumerate}
    \item \textit{Linearity:} For $f, g \in L^1(\R^d)$ and constants $a, b \in \R$,
    \begin{equation*}
        \mathcal{F}[af + bg] = a\mathcal{F}[f] + b\mathcal{F}[g].
    \end{equation*}
    \item \textit{Space Shifting and Scaling:} Let $\v{A} \in \R^{d \times d}$ be an invertible matrix and $\v{b} \in \R^d$ a constant vector. Let $f \in L^1(\R^d)$ and define $h(x) \coloneqq f(\v{Ax} + \v{b})$. Then
    \begin{equation*}
        \mathcal{F}[h](t) = \e^{\ii 2\pi (\v{A}^{-1} t)^\T \v{b}} \widetilde{f}(\v{A}^{-\T} t) / |\det(\v{A})|,
    \end{equation*}
    where $\v{A}^{-\T} \coloneqq (\v{A}^\T)^{-1} = (\v{A}^{-1})^\T$.
    \item \textit{Frequency Shifting and Scaling:} Let $\v{A} \in \R^{d \times d}$ be an invertible matrix and $\v{b} \in \R^d$ a constant vector. Let $f \in L^1(\R^d)$ and define
    \begin{equation*}
        h(x) \coloneqq \e^{-\ii 2\pi \v{b}^\T \v{A}^{-\T} \v{x}} f(\v{A}^{-\T} \v{x}) / |\det(\v{A})|.
    \end{equation*}
    Then $\mathcal{F}[h](t) = \widetilde{f}(\v{A}t + \v{b})$.
    \item \textit{Differentiation:} Let $f \in L^1(\R^d) \cap C^1(\R^d)$ and let $f_k \coloneqq \partial f / \partial x_k$ be the partial derivative of $f$ with respect to $x_k$. If $f_k \in L^1(\R^d)$ for $k=1, \dots, d$, then
    \begin{equation*}
        \mathcal{F}[f_k](t) = (\ii 2\pi t_k) \widetilde{f}(t).
    \end{equation*}
    \item \textit{Convolution:} Let $f, g \in L^1(\R^d)$ be real or complex valued functions. Their convolution, $f * g$, is defined as
    \begin{equation*}
        (f * g)(x) = \int_{\R^d} f(y) g(x-y) \di y,
    \end{equation*}
    and is also in $L^1(\R^d)$. Moreover, the Fourier transform satisfies
    \begin{equation*}
        \mathcal{F}[f * g] = \mathcal{F}[f]\mathcal{F}[g].
    \end{equation*}
    \item \textit{Duality:} Let $f$ and $\mathcal{F}[f]$ both be in $L^1(\R^d)$. Then $\mathcal{F}[\mathcal{F}[f]](t) = f(-t)$.
    \item \textit{Product Formula:} Let $f, g \in L^1(\R^d)$ and denote by $\widetilde{f}, \widetilde{g}$ their respective Fourier transforms. Then $\widetilde{f}g, f\widetilde{g} \in L^1(\R^d)$, and
    \begin{equation*}
        \int_{\R^d} \widetilde{f}(z) g(z) \di z = \int_{\R^d} f(z) \widetilde{g}(z) \di z.
    \end{equation*}
\end{enumerate}

There are many additional properties which hold if $f \in L^1(\R^d) \cap L^2(\R^d)$. In particular, if $f, g \in L^1(\R^d) \cap L^2(\R^d)$, then $\widetilde{f}, \widetilde{g} \in L^2(\R^d)$ and $\langle \widetilde{f}, \widetilde{g} \rangle = \langle f, g \rangle$, a result often known as \textit{Parseval's formula}. Putting $g=f$ gives the result often referred to as \textit{Plancherel's theorem}.

The Fourier transform can be extended in several ways, in the first instance to functions in $L^2(\R^d)$ by continuity. A substantial extension of the theory is realized by replacing integration with respect to the Lebesgue measure (i.e., $\int_{\R^d} \cdots \di x$) with integration with respect to a (finite Borel) measure $\mu$ (i.e., $\int_{\R^d} \cdots \mu(\di x)$). Moreover, there is a close connection between the Fourier transform and characteristic functions arising in probability theory. Indeed, if $\v{X}$ is a random vector with pdf $f$, then its characteristic function $\psi$ satisfies
\begin{equation*}
    \psi(t) \coloneqq \E \e^{\ii t^\T \v{X}} = \mathcal{F}[f](-t/(2\pi)).
\end{equation*}

\subsubsection{Discrete Fourier Transform}
Here, we introduce the (univariate) discrete Fourier transform, which can be viewed as a special case of the Fourier transform introduced in Definition A.6, where $d=1$, integration is with respect to the counting measure, and $f(x) = 0$ for $x < 0$ and $x > (n-1)$.

\begin{defn}{Discrete Fourier Transform}{DiscreteFourierTransform}
The \textit{discrete Fourier transform} (DFT) of a vector $\v{x} = [x_0, \dots, x_{n-1}]^\T \in \C^n$ is the vector $\widetilde{\v{x}} = [\widetilde{x}_0, \dots, \widetilde{x}_{n-1}]^\T$ whose elements are given by
\begin{equation}
    \widetilde{x}_t = \sum_{s=0}^{n-1} \omega^{st} x_s, \quad t = 0, \dots, n-1, \label{eq:DFTDefinition}
\end{equation}
where $\omega = \exp(-\ii 2\pi/n)$.
\end{defn}

In other words, $\widetilde{\v{x}}$ is obtained from $\v{x}$ via the linear transformation
\begin{equation*}
    \widetilde{\v{x}} = \v{F} \v{x},
\end{equation*}
where
\begin{equation*}
    \v{F} =
    \begin{bmatrix}
        1 & 1 & 1 & \dots & 1 \\
        1 & \omega & \omega^2 & \dots & \omega^{n-1} \\
        1 & \omega^2 & \omega^4 & \dots & \omega^{2(n-1)} \\
        \vdots & \vdots & \vdots & \ddots & \vdots \\
        1 & \omega^{n-1} & \omega^{2(n-1)} & \dots & \omega^{(n-1)^2}
    \end{bmatrix}.
\end{equation*}
The matrix $\v{F}$ is a so-called \textit{Vandermonde matrix}, and is clearly symmetric (i.e., $\v{F} = \v{F}^\T$). Moreover, $\v{F}/\sqrt{n}$ is in fact a unitary matrix and hence its inverse is simply its complex conjugate $\overline{\v{F}}/\sqrt{n}$. Thus, $\v{F}^{-1} = \overline{\v{F}}/n$ and we have that the \textit{inverse discrete Fourier transform} (IDFT) is given by
\begin{equation}
    x_t = \frac{1}{n} \sum_{s=0}^{n-1} \omega^{-st} \widetilde{x}_s, \quad t = 0, \dots, n-1, \label{eq:IDFTDefinition}
\end{equation}
or in terms of matrices and vectors,
\begin{equation*}
    \v{x} = \overline{\v{F}} \widetilde{\v{x}} / n.
\end{equation*}
Observe that the IDFT of a vector $\v{y}$ is related to the DFT of its complex conjugate $\overline{\v{y}}$, since
\begin{equation*}
    \overline{\v{F}} \v{y} / n = \overline{\v{F} \overline{\v{y}}} / n.
\end{equation*}
Consequently, an IDFT can be computed via a DFT.
There is a close connection between circulant matrices $\v{C}$ and the DFT. To make this connection concrete, let $\v{C}$ be the circulant matrix corresponding to the vector $\v{c} \in \C^n$ and denote by $\v{f}_t$ the $t$-th column of the discrete Fourier matrix $\v{F}$, $t=0, 1, \dots, n-1$. Then, the $s$-th element of $\v{C} \v{f}_t$ is
\begin{equation*}
    \sum_{k=0}^{n-1} c_{(s-k) \bmod n} \omega^{tk} = \sum_{y=0}^{n-1} c_y \omega^{t(s-y)} = \underbrace{\omega^{ts}}_{s\text{-th element of } \v{f}_t} \underbrace{\sum_{y=0}^{n-1} c_y \omega^{-ty}}_{\lambda_t}.
\end{equation*}
Hence, the eigenvalues of $\v{C}$ are
\begin{equation*}
    \lambda_t = \v{c}^\T \overline{\v{f}}_t, \quad t = 0, 1, \dots, n-1,
\end{equation*}
with corresponding eigenvectors $\v{f}_t$. Collecting the eigenvalues into the vector $\blambda = [\lambda_0, \dots, \lambda_{n-1}]^\T = \overline{\v{F}} \v{c}$, we therefore have the eigen-decomposition
\begin{equation*}
    \v{C} = \v{F} \diag(\blambda) \overline{\v{F}} / n.
\end{equation*}
Consequently, one can compute the \textit{circular convolution} of a vector $\v{a} = [a_1, \dots, a_n]^\T$ and $\v{c} = [c_0, \dots, c_{n-1}]^\T$ by a series of DFTs as follows. Construct the circulant matrix $\v{C}$ corresponding to $\v{c}$. Then, the circular convolution of $\v{a}$ and $\v{c}$ is given by $\v{y} = \v{Ca}$. Proceed in four steps:
\begin{enumerate}
    \item Compute $\v{z} = \overline{\v{F}} \v{a} / n$.
    \item Compute $\blambda = \overline{\v{F}} \v{c}$.
    \item Compute $\v{p} = \v{z} \odot \blambda = [z_1 \lambda_0, \dots, z_n \lambda_{n-1}]^\T$.
    \item Compute $\v{y} = \v{Fp}$.
\end{enumerate}
Steps 1 and 2 are (up to constants) in the form of an IDFT, and step 4 is in the form of a DFT. These are computable via the FFT in $\bigO(n \ln n)$ time. Step 3 is a dot product computable in $\bigO(n)$ time. Thus, the circular convolution can be computed with aid of the FFT in $\bigO(n \ln n)$ time.

One can also efficiently compute the product of an $n \times n$ Toeplitz matrix $\v{T}$ and an $n \times 1$ vector $\v{a}$ in $\bigO(n \ln n)$ time by embedding $\v{T}$ into a circulant matrix $\v{C}$ of size $2n \times 2n$. Namely, define
\begin{equation*}
    \v{C} =
    \begin{bmatrix}
        \v{T} & \v{B} \\
        \v{B} & \v{T}
    \end{bmatrix},
\end{equation*}
where
\begin{equation*}
    \v{B} =
    \begin{bmatrix}
        0 & t_{n-1} & \dots & t_2 & t_1 \\
        t_{-(n-1)} & 0 & t_{n-1} & \dots & t_2 \\
        \vdots & t_{-(n-1)} & 0 & \ddots & \vdots \\
        t_{-2} & \ddots & \ddots & \ddots & t_{n-1} \\
        t_{-1} & t_{-2} & \dots & t_{-(n-1)} & 0
    \end{bmatrix}.
\end{equation*}
Then a product of the form $\v{y} = \v{Ta}$ can be computed in $\bigO(n \ln n)$ time, since we may write
\begin{equation*}
    \v{C}
    \begin{bmatrix}
        \v{a} \\ \v{0}
    \end{bmatrix}
    =
    \begin{bmatrix}
        \v{T} & \v{B} \\
        \v{B} & \v{T}
    \end{bmatrix}
    \begin{bmatrix}
        \v{a} \\ \v{0}
    \end{bmatrix}
    =
    \begin{bmatrix}
        \v{Ta} \\ \v{Ba}
    \end{bmatrix}.
\end{equation*}
The left-hand side is a product of a $2n \times 2n$ circulant matrix with vector of length $2n$, and so can be computed in $\bigO(n \ln n)$ time via the FFT, as previously discussed.

Conceptually, one can also solve equations of the form $\v{Cx} = \v{b}$ for a given vector $\v{b} \in \C^n$ and circulant matrix $\v{C}$ (corresponding to $\v{c} \in \C^n$, assuming all its eigenvalues are non-zero) via the following four steps:
\begin{enumerate}
    \item Compute $\v{z} = \overline{\v{F}} \v{b} / n$.
    \item Compute $\blambda = \overline{\v{F}} \v{c}$.
    \item Compute $\v{p} = \v{z} / \blambda = [z_1 / \lambda_0, \dots, z_n / \lambda_{n-1}]^\T$.
    \item Compute $\v{x} = \v{Fp}$.
\end{enumerate}
Once again, Steps 1 and 2 are (up to constants) in the form of an IDFT, and Step 4 is in the form of a DFT, all of which are computable via the FFT in $\bigO(n \ln n)$ time, and Step 3 is computable in $\bigO(n)$ time, meaning the solution $\v{x}$ can be computed using the FFT in $\bigO(n \ln n)$ time.

\subsubsection{Fast Fourier Transform}
The \textit{fast Fourier transform} (FFT) is a numerical algorithm for the fast evaluation of \eqref{eq:DFTDefinition} and \eqref{eq:IDFTDefinition}. By using a divide-and-conquer strategy, the algorithm reduces the computational complexity from $\bigO(n^2)$ (for the naïve evaluation of the linear transformation) to $\bigO(n \ln n)$.

The essence of the algorithm lies in the following observation. Suppose $n = r_1 r_2$. Then one can express any index $t$ appearing in \eqref{eq:DFTDefinition} via a pair $(t_0, t_1)$, with $t = t_1 r_1 + t_0$, where $t_0 \in \{0, 1, \dots, r_1 - 1\}$ and $t_1 \in \{0, 1, \dots, r_2 - 1\}$. Similarly, one can express any index $s$ appearing in \eqref{eq:DFTDefinition} via a pair $(s_0, s_1)$, with $s = s_1 r_2 + s_0$, where $s_0 \in \{0, 1, \dots, r_2 - 1\}$ and $s_1 \in \{0, 1, \dots, r_1 - 1\}$.

Identifying $\widetilde{x}_t \equiv \widetilde{x}_{t_1, t_0}$ and $x_s \equiv x_{s_1, s_0}$, we may re-express \eqref{eq:DFTDefinition} as
\begin{equation}
    \widetilde{x}_{t_1, t_0} = \sum_{s_0=0}^{r_2-1} \omega^{s_0 t} \sum_{s_1=0}^{r_1-1} \omega^{s_1 r_2 t} x_{s_1, s_0}, \quad t_0 = 0, 1, \dots, r_1 - 1, t_1 = 0, 1, \dots, r_2 - 1. \label{eq:FFTStep}
\end{equation}
Observe that $\omega^{s_1 r_2 t} = \omega^{s_1 r_2 t_0}$ (because $\omega^{r_1 r_2} = 1$), so that the inner sum over $s_1$ depends only on $s_0$ and $t_0$. Define
\begin{equation*}
    y_{t_0, s_0} \coloneqq \sum_{s_1=0}^{r_1-1} \omega^{s_1 r_2 t_0} x_{s_1, s_0}, \quad t_0 = 0, 1, \dots, r_1 - 1, s_0 = 0, 1, \dots, r_2 - 1.
\end{equation*}
Computing each $y_{t_0, s_0}$ requires $\bigO(n r_1)$ operations. In terms of the $\{y_{t_0, s_0}\}$, \eqref{eq:FFTStep} can be written as
\begin{equation*}
    \widetilde{x}_{t_1, t_0} = \sum_{s_0=0}^{r_2-1} \omega^{s_0 t} y_{t_0, s_0}, \quad t_1 = 0, 1, \dots, r_2 - 1, t_0 = 0, 1, \dots, r_1 - 1,
\end{equation*}
requiring $\bigO(n r_2)$ operations to compute. Thus, calculating the DFT using this two-step procedure requires $\bigO(n(r_1 + r_2))$ operations, rather than $\bigO(n^2)$.

Now supposing $n = r_1 r_2 \cdots r_m$, repeated application the above divide-and-conquer idea yields an $m$-step procedure requiring $\bigO(n(r_1 + r_2 + \dots + r_m))$ operations. In particular, if $r_k = r$ for all $k=1, 2, \dots, m$, we have that $n = r^m$ and $m = \log_r n$, so that the total number of operations is $\bigO(r n m) \equiv \bigO(r n \log_r(n))$. Typically, the radix $r$ is a small (not necessarily prime) number, for instance $r=2$.

\section{Multivariate Differentiation and Optimisation}
The purpose of this appendix is to review various aspects of multivariate differentiation and optimisation. We assume the reader is familiar with differentiating a real-valued function.